\section*{第 \ref{ch:b3:dna-repair} 章——基因组稳定性：DNA 损伤、修复与重组}

\begin{solution}{exo:b3:dna-repair:1}
尿嘧啶：自发产生（C 的\omterm{def:b3:dna-repair:lesions}{脱氨基}），\omterm{def:b3:dna-repair:excision}{碱基切除修复}由
尿嘧啶-DNA 糖基化酶完成。胸腺\omterm{def:b3:dna-repair:lesions}{嘧啶二聚体}：环境因素（紫外线）引起，由\omterm{def:b3:dna-repair:excision}{核苷酸
切除修复}完成。\omterm{def:b3:dna-repair:lesions}{8-氧代鸟嘌呤}：自发产生（氧化），\omterm{def:b3:dna-repair:excision}{碱基切除
修复}，由 8-oxoG 糖基化酶完成。复制叉后方的 G--T 错配：
复制错误，\omterm{prop:b3:dna-repair:fidelity}{错配修复}。G1 期的\omterm{def:b3:dna-repair:lesions}{双链断裂}：
环境因素（辐射）引起或自发产生，\omterm{def:b3:dna-repair:dsb}{非同源末端连接}，
因无可用的姐妹染色单体。
\end{solution}

\begin{solution}{exo:b3:dna-repair:2}
核苷酸切除识别所有大型\omterm{def:b3:dna-repair:lesions}{损伤}共有的一种物理性质
——螺旋的扭曲，或 RNA 聚合酶停滞——
并切除含有致损原因的一段链。小\omterm{def:b3:dna-repair:lesions}{损伤}
（尿嘧啶、8-oxoG、烷基化碱基）几乎不扭曲螺旋，因此它们
必须依其化学性质加以识别，每种错误碱基
各有一种糖基化酶。
\end{solution}

\begin{solution}{exo:b3:dna-repair:3}
错配由两个正常碱基构成，而只有新链上的那个碱基是错的；
因此系统必须能分辨哪一条链是新合成的。
\emph{E.~coli} 用甲基在 GATC 位点标记亲本链，
新链在几分钟内缺少甲基，MutH 切开
未甲基化的链。在 \emph{dam} 突变体中，两条链均未甲基化，
MutH 切开哪条链都可以，修复时去除正确碱基与错误碱基的机会相当，
使一半错误固定为突变（即突变子）——
两条链上邻近 GATC 位点处的切口
形成\omterm{def:b3:dna-repair:lesions}{双链断裂}。
\end{solution}

\begin{solution}{exo:b3:dna-repair:4}
约 $35\times 2 = 70$ 个\omterm{def:b3:dna-repair:lesions}{双链断裂}。一小时后大部分已
被连接：半衰期约为 \qty{30}{min}，约 $70/4 \approx 18$ 个
焦点。六小时后几乎全部被修复（$70/2^{12} < 1$），
只留下少数持续存在的复杂断裂。
\end{solution}

\begin{solution}{exo:b3:dna-repair:5}
(a) $10^{-5}\times 10^{-2}\times 10^{-3} = 10^{-10}$；$6.4\times
10^{9}\times 10^{-10} = 0.64$，即每次分裂 0.64 个突变。(b) 无\omterm{prop:b3:dna-repair:fidelity}{错配修复}时
为 $10^{-7}$，每次分裂 $640$ 个。(c) 无\omterm{prop:b3:dna-repair:fidelity}{校对}功能时
为 $10^{-8}$，每次分裂 $64$ 个。\omterm{prop:b3:dna-repair:fidelity}{错配修复}缺陷的突变体
更危险，其突变率比正常高出千倍。
\end{solution}

\begin{solution}{exo:b3:dna-repair:6}
$k = \ln 2/\qty{0.25}{h} = \qty{2.8}{h^{-1}}$; $D = 10^{4}/24 =
\qty{420}{h^{-1}}$；$L^{*} = 420/2.8 \approx 150$ 个无碱基位点同时存在
任一时刻。修复速度慢十倍时：$L^{*} \approx \num{1500}$。复制叉
每个基因座恰好经过一次，遇到该处此刻存在的\omterm{def:b3:dna-repair:lesions}{损伤}；
对全基因组取平均，它平均遇到约 $L^{*}$ 个
——正常时约 $150$ 个，加药后为 $\num{1500}$ 个。
\end{solution}

\begin{solution}{exo:b3:dna-repair:7}
$k_{\text{NHEJ}} = \ln 2/0.5 = \qty{1.39}{h^{-1}}$, $k_{\text{HR}} =
\ln 2/4 = \qty{0.17}{h^{-1}}$；HR 所占比例 $0.17/1.56 = 0.11$；
半衰期 $\ln 2/1.56 = \qty{27}{min}$。塌陷的复制叉产生一个
\emph{单端}断裂：没有第二个末端供 NHEJ 连接，因此
只有与姐妹染色单体重组才能恢复复制叉——况且这
在任何情况下都是唯一准确的途径。
\end{solution}

\begin{solution}{exo:b3:dna-repair:8}
(a) \omterm{def:b3:dna-repair:mmr}{林奇综合征}：微卫星不稳定（肿瘤
细胞之间长度不一），结直肠癌、子宫内膜癌、胃癌和卵巢癌，日光反应
正常。(b) XP-A：微卫星稳定，微量日晒即
严重晒伤，儿童期即患皮肤癌和眼癌，并且
该组患者还伴有进行性神经退化。(c) XP-V：
微卫星稳定，切除修复正常，因此晒伤较轻，
但未被切除的二聚体由易错
聚合酶绕过——导致皮肤癌，发病年龄晚于 A 组。
\end{solution}

\begin{solution}{exo:b3:dna-repair:9}
肝细胞：Cre 已从两个等位基因上切除外显子 3（以环状形式
丢失），因此基因从出生起在每个肝细胞中即失活
——肝脏特异性敲除。神经元：无 Cre，两侧带 loxP 的等位基因
完整且有功能。反向排列：Cre 使外显子
在两种方向之间不断翻转，因此任一时刻约一半的
细胞携带倒位而无功能的外显子——形成镶嵌体，而不是干净的
缺失。
\end{solution}

\begin{solution}{exo:b3:dna-repair:10}
随着 \omterm{prop:b3:dna-repair:sos}{LexA} 被切割，其浓度逐渐下降。结合 \omterm{prop:b3:dna-repair:sos}{LexA} 较弱的
操纵基因在浓度稍降时即被释放——\emph{uvrA}，
\emph{uvrB}、\emph{recA}——因此准确修复率先启动；
\emph{sulA} 及易错聚合酶的操纵基因与 \omterm{prop:b3:dna-repair:sos}{LexA} 结合
紧密，只有当 \omterm{prop:b3:dna-repair:sos}{LexA} 几乎耗尽时才被释放，也就是
\omterm{def:b3:dna-repair:lesions}{损伤}持续很长时间之后。不含易错聚合酶的细菌
会在被其无法绕过的\omterm{def:b3:dna-repair:lesions}{损伤}所阻断的复制叉处死亡，
并且由于没有诱导突变，对抗生素产生抗性的
进化会慢得多——对细菌不利，
却对患者有利；因此 SOS 反应的抑制剂
正是出于这一原因而受到研究。
\end{solution}

\begin{solution}{exo:b3:dna-repair:11}
链式反应：PARP 被抑制 $\to$ 单链断裂持续存在 $\to$ 一个
复制叉将每个断裂转化为单端\omterm{def:b3:dna-repair:lesions}{双链断裂}
$\to$ 修复需要\omterm{def:b3:dna-repair:dsb}{同源重组} $\to$ 肿瘤细胞因缺乏 \omterm{def:b3:dna-repair:dsb}{BRCA}，
无法 $\to$ 错误连接、染色体丢失、死亡。正常
细胞具有一个有功能的等位基因，可以重组并存活。抗性：
（1）\emph{\omterm{def:b3:dna-repair:dsb}{BRCA}} 上的第二次突变恢复了阅读框
并恢复了蛋白质功能——可通过对肿瘤或
血液中循环的肿瘤 DNA 测序检测，并通过照射后 \omterm{def:b3:dna-repair:dsb}{Rad51} 焦点的重现检测；
（2）阻断末端切除的蛋白丢失
（53BP1--shieldin），可在没有 BRCA1 时部分恢复重组
——可通过染色时这些蛋白不再可见及 \omterm{def:b3:dna-repair:dsb}{Rad51} 焦点恢复来检测；
此外，阻止 PARP1 滞留于 DNA 上的突变，或药物
外排泵，分别通过测序和表达检测。
\end{solution}

\begin{solution}{exo:b3:dna-repair:12}
两条同源双链体，一条带 $A$，另一条带 $a$，在标记两侧交换
单链，使该区段内每条双链各含一条来自各亲本的链，形成
一个带有 $A/a$ 错配的异源双链体。若
\omterm{prop:b3:dna-repair:fidelity}{错配修复}将入侵双链体上的错配改为 $A$，
三条染色单体带 $A$、一条带 $a$：$3{:}1$；改为 $a$ 时为
$1{:}3$；若未修复，染色单体在减数分裂后的第一次有丝分裂中分离出 $A$ 和 $a$，
形成两种不同的子囊孢子
（减数分裂后分离，八孢子子囊中比例为 $5{:}3$）。
异源双链区段仅为起始断裂周围的数百至数千个碱基对，
因此只有其中的标记能被转换，
而它们按构造必然紧邻交换点。
\end{solution}

\begin{solution}{pb:b3:dna-repair:1}
\textbf{1.} $10^{4} + 400 + 10^{4} \approx 2\times 10^{4}$ 个\omterm{def:b3:dna-repair:lesions}{损伤}，每天
即每年 $7\times 10^{6}$ 个——约每年每千碱基基因组一个
\omterm{def:b3:dna-repair:lesions}{损伤}。
\textbf{2.} $10^{-5}\times 10^{-2}\times 10^{-3} = 10^{-10}$ 每碱基对；
$6.4\times 10^{9}\times 10^{-10} = 0.64$，即每次分裂的突变数。
\textbf{3.} $0.64\times 0.015 \approx 0.01$ 个编码突变，每
次分裂；经过 $50$ 次分裂，约为 $0.5$。
\textbf{4.} 每碱基对 $10^{-7}$，每次分裂 $640$ 个突变。发生滑脱的
细胞：无修复时 $10^{-4}\times 40\times 10^{9} = 4\times 10^{6}$；
有修复时 $10^{-7}\times 40\times 10^{9} = 4000$——相差千倍，
表现为\omterm{def:b3:dna-repair:mmr}{微卫星不稳定性}。
\textbf{5.} 尿嘧啶不属于 DNA，会被尿嘧啶-DNA
糖基化酶识别并切除，得以修复。由 5-甲基胞嘧啶产生的胸腺嘧啶是
与 G 配对的正常碱基；复制之外没有任何标记表明它有误，
于是变为 C$\to$T 突变——这正是使基因组中
\omterm{def:b3:chromatin-epigenetics:methylation}{CpG} 减少的机制。
\textbf{6.} 8-oxoG 糖基化酶在被氧化的碱基
仍与 C 配对时（复制之前）将其切除；第二种糖基化酶
切除与 8-oxoG 对位错掺入的 A（复制之后，使 G 有
再次修复的机会）；水解酶破坏被氧化的 dGTP，使其
不会在 A 对位被掺入。若所有修复均失效，8-oxoG 与 A 配对，
下一轮复制便固定为 G$\cdot$C$\to$T$\cdot$A 颠换。
\textbf{7.} 正常细胞 $k = \ln 2/2 = \qty{0.35}{h^{-1}}$；XP $k = \ln 2/70
= \qty{0.0099}{h^{-1}}$。
\textbf{8.} 正常 $10^{5} e^{-2.77} \approx 6200$；XP $10^{5}
e^{-0.079} \approx \num{92000}$。
\textbf{9.} 每个细胞的突变数为 $6.2$ 和 $92$；编码突变数为 $0.09$ 和 $1.4$。
\textbf{10.} $500$ 次照射产生 $47$ 个编码突变（正常），
$690$ 个（XP），相比复制带来的 $0.5$——阳光主导了
即便是正常人，皮肤的突变负荷也主要来自阳光。
\textbf{11.} 编码序列长 $9.6\times 10^{7}$ 碱基对；驱动基因
占比 $4\times 10^{4}/9.6\times 10^{7} = 4.2\times 10^{-4}$。平均
命中数为 $690\times 4.2\times 10^{-4} = 0.29$；$P(\ge 1) = 1 - e^{-0.29} =
0.25$。在 $10^{9}$ 个细胞中，$2.5\times 10^{8}$ 个携带驱动突变
（正常儿童：平均 $0.02$，\qty{2}{\%}）。
\textbf{12.} 二聚体本身不是突变；只有当复制叉以易错旁路方式复制它时，才会成为突变，
因此积累大量未修复二聚体的分裂细胞
最易发生突变。紫外线只能到达皮肤和眼睛，
内部组织几乎不产生二聚体，因而在
XP 患者中几乎正常。
\textbf{13.} $70$ 个断裂；仅依靠 NHEJ 时经过 \qty{1}{h} 后，$70\times
2^{-2} \approx 18$ 个焦点。
\textbf{14.} $k_{\text{NHEJ}} = \qty{1.39}{h^{-1}}$，$k_{\text{HR}} =
\qty{0.17}{h^{-1}}$；HR 所占比例 $0.17/1.56 = 0.11$。
\textbf{15.} 半衰期 $\ln 2/1.56 = \qty{27}{min}$；经过 \qty{1}{h} 后
$70\,e^{-1.56} \approx 15$ 个断裂。
\textbf{16.} 每个细胞被破坏的基因数约为 $70\times 0.015\times 0.7 \approx 0.7$
个。
\textbf{17.} $n = 70$：$2415$ 对 $\times 3\times 10^{-4} = 0.72$
次易位。$n = 7$：$21\times 3\times 10^{-4} = 0.006$。断裂
数 $\propto D$，而同时发生的断裂\emph{对}数
则 $\propto n^{2} \propto D^{2}$。
\textbf{18.} $D = \qty{7}{h^{-1}}$，$L^{*} = 7/1.39 \approx 5$ 个断裂
同时存在：$10$ 对，$0.003$ 次易位——比同剂量
瞬时照射少两百倍。分次照射使正常组织在各次之间
得以修复，并使同时存在的断裂数保持较少。
\textbf{19.} 修复完好时：$e^{-2/1.5} = 0.26$（\qty{2}{Gy} 时）；
$e^{-4} = 0.018$（\qty{6}{Gy} 时）。无 NHEJ 时为 $e^{-4} = 0.018$ 和
$e^{-12} = 6\times 10^{-6}$。$D_{0}$ 是平均每个细胞留下
一个致死性未修复\omterm{def:b3:dna-repair:lesions}{损伤}（存活率为 $1/e$）的剂量。
\textbf{20.} 本可由 HR 准确修复的两端 G2 期断裂中的 \qty{11}{\%}，
现在改由 NHEJ 处理，而每一个单端复制叉断裂——NHEJ 根本无法
正确修复它——都会被错误连接：经年累月，基因组中
积累缺失、易位及拷贝数变化，
这些正是 \emph{\omterm{def:b3:dna-repair:dsb}{BRCA}} 肿瘤的特征性痕迹。
\textbf{21.} 每天 $10^{4}\times 0.01 = 100$ 个单端断裂。一个
单端断裂没有配对的末端；NHEJ 只能让它保持原样，或将其
与某个无关的末端连接，从而导致易位或缺失。
\textbf{22.} 一个功能正常的 \emph{BRCA2} 等位基因即可产生足以用于重组的蛋白质：
正常细胞修复其复制叉断裂，
并能在抑制剂作用下存活。
\textbf{23.} 第一个移码下游的第二个移码恢复了
阅读框，产生缩短但部分有功能的 BRCA2，因此
重组恢复，药物不再致死。可通过对耐药肿瘤或
循环肿瘤 DNA 中的 \emph{BRCA2} 测序检测，
并通过照射后 \omterm{def:b3:dna-repair:dsb}{Rad51} 焦点的重现
进行功能检测。
\textbf{24.} 杀伤作用需要\omterm{prop:b3:dna-repair:fidelity}{错配修复}去识别被烷基化的
碱基；缺乏该修复的\omterm{def:b3:dna-repair:mmr}{林奇综合征}肿瘤具有耐受性——药物无法杀死
它（此类肿瘤改由其他手段治疗，包括
免疫疗法，因其突变负荷很高而对其
敏感）。
\textbf{25.} XP 细胞中复制叉处约有 \num{92000} 个二聚体；
每碱基对每次分裂的突变率为 $10^{-10}$；HR 修复
\qty{11}{\%} 的 G2 期断裂。
\end{solution}
